\documentclass{report}
\usepackage{import}
\usepackage{tcolorbox}
\usepackage{amsmath}
\usepackage{amssymb}
\usepackage{amsthm}
\usepackage{geometry}
\usepackage{enumitem}
\usepackage{pdfpages}
\usepackage{transparent}
\usepackage[hidelinks]{hyperref}
\usepackage{booktabs}
\usepackage{bm}
\usepackage{tabularx}
\tcbuselibrary{theorems}
\tcbuselibrary{skins}

% tabularx Settings
\renewcommand\tabularxcolumn[1]{>{\centering\arraybackslash}m{#1}}

% Page Settings
\newgeometry{vmargin={15mm}, hmargin={15mm, 15mm}}
\setlength{\parindent}{0pt}
\setlength{\parskip}{1em plus 1pt minus 1pt}

% hyperref Settings
\hypersetup{
    colorlinks=true,
    linkcolor=blue,
    filecolor=magenta,      
    urlcolor=cyan,
}

% inkscape Settings
\newcommand{\incfig}[1]{%
    \def\svgwidth{\columnwidth}
    \import{./figure/}{#1.pdf_tex}
}

% Definition Box Styles
\newtcbtheorem[auto counter, number within=subsection]{defn}{Definition}{
% basic settings
    theorem name and number,
    separator sign={},
    description delimiters={(}{)},
    enhanced jigsaw,
    sharp corners,
% box margins and rules
    boxrule=1pt,
    left=0.2cm,right=0.2cm,top=0.2cm,
    toptitle=0.1cm,
    bottomtitle=0.08cm,
% box colors
    colframe=white!25!black,
    colback=white,
    coltitle=white,
% box fonts
    fonttitle=\bfseries,fontupper=\normalsize
}{defn}

% Theorem Box Styles
\newtcbtheorem[auto counter, number within=subsection]{thm}{Theorem}{
% basic settings
    theorem name and number,
    separator sign={},
    description delimiters={(}{)},
    enhanced jigsaw,
    sharp corners,
% box margins and rules
    boxrule=1pt,
    left=0.2cm,right=0.2cm,top=0.2cm,
    toptitle=0.1cm,
    bottomtitle=0.08cm,
% box colors
    colframe=white!25!black,
    colback=white,
    coltitle=white,
% box fonts
    fonttitle=\bfseries,fontupper=\normalsize
}{thm}

% Lemma Box Styles
\newtcbtheorem[number within=tcb@cnt@thm]{lem}{Lemma}{
% basic settings
    theorem name and number,
    separator sign={},
    description delimiters={(}{)},
    enhanced jigsaw,
    sharp corners,
% box margins and rules
    boxrule=1pt,
    left=0.2cm,right=0.2cm,top=0.2cm,
    toptitle=0.1cm,
    bottomtitle=0.08cm,
% box colors
    colframe=white!50!black,
    colback=white,
    coltitle=white,
% box fonts
    fonttitle=\bfseries,fontupper=\normalsize
}{lem}

% Corollary Box Styles
\newtcbtheorem[number within=tcb@cnt@thm]{cor}{Corollary}{
% basic settings    
    theorem name and number,
    separator sign={},
    description delimiters={(}{)},
    enhanced jigsaw,
    sharp corners,
% box margins and rules
    boxrule=1pt,
    left=0.2cm,right=0.2cm,top=0.2cm,
    toptitle=0.1cm,
    bottomtitle=0.08cm,
% box colors
    colframe=white!50!black,
    colback=white,
    coltitle=white,
% box fonts
    fonttitle=\bfseries,fontupper=\normalsize
}{cor}

% Remark Box Styles
\newtcbtheorem[number within=subsection]{rem}{Remark}{
% basic settings
    theorem name and number,
    separator sign={},
    description delimiters={(}{)},
    enhanced jigsaw,
    sharp corners,
% box colors
    colframe=white!60!black,
    colback=white,
    coltitle=black,
% box margins and rules
    boxrule=1pt,
    left=0.2cm,right=0.2cm,top=0.2cm,
    toptitle=0.1cm,
    bottomtitle=0.08cm,
% box fonts
    fonttitle=\bfseries,fontupper=\normalsize
}{rem}

% Proof Box Styles

\title{Notes on Models of Machine Learning}
\author{Hyoung Joon, Kim}
\date{\today}


\begin{document}

\maketitle
\tableofcontents
\newpage

\chapter{Introduction}
This is a set of notes about machine learning, based on the book by Bishop and Murphy.
I assume that the reader is familiar with ML topics, since this is just a simple reference.

    \section{Notations}
        \begin{itemize}
            \item real value: $x$
            \item vector value: $\bs{x}$
            \item vector: $\bs{x} = (x_1, x_2, \ldots, x_n)
            = 
            \begin{pmatrix}
                x_1 \\
                x_2 \\
                \vdots \\
                x_n    
            \end{pmatrix}
            =
            (x_1\ x_2\ \ldots\ x_n)^\top
            $
            \item matrix: $\textbf{X}$
        \end{itemize}

\chapter{Linear Regression Models}
    In this chapter, we discuss linear regression, which is a widely used method for predicting a real-valued output.
    The most noticable property of linear regression model is that the expected value of the model is linear in parameters,
    i.e., $\mathbb{E}[y|\bs{x}]=\bs{w}^{\top}\bs{x}$.

    \section{Least Squares Linear Regression}
        \subsection{Model Explanation}
            \begin{defn}{Least Squares Linear Regression Model}{LSLR}
                The following is the most common form of the linear regression model:
                \[
                p(y|\bs{x}, \param) = \mathcal{N}(y|w_0+\bs{w}^\top \bs{x}, \sigma^2)
                \]
                where $w_0$ is called \textbf{bias} and $\bs{w}$ is called \textbf{weights}. If the output is an N-dimensional vector, then the model has following form:
                \[ 
                p(\bs{y}|\bs{x}, \param) = \prod_{k=1}^{N}\mathcal{N}(y_k|\bs{w}_k^\top\bs{x}, \sigma_{k}^{2})
                \]
                (the model is assuming that each elements of the output vector is independent.)
            \end{defn}

            We will usually assume that $\bs{x}$ is written as $(1, x_1, x_2, \ldots, x_n)$ so that the bias $w_0$ can be absorbed into the weight vector $\bs{w}$.

        \subsection{Parameter Estimation}
            The estimation of parameters of \hyperref[defn:LSLR]{this model} is usually done through calculating the MLE.

            \begin{thm}{MLE of Least Squares Linear Regression}{MLEofLSLR}
                Let $\data=\{(\bs{x}_k, y_k)|\ k=1,2,\ldots,N\}$ be the given data, and $\hat{y_k}=\bs{w}^\top \bs{x}_k$ the predicted data (using the plug-in predictive distribution).
                Let $\textbf{X}=(\bs{x}_1 \ \bs{x}_2 \ \ldots \ \bs{x}_N)^\top$ and $\textbf{y}=(y_1,y_2,\ldots,y_N)$
                \begin{itemize}
                    \item objective function: NLL (negative log likelihood) = $-\sum_{k=1}^{N}\log{\mathcal N(y_k|\bs{w}^\top \bs{x}_k, \sigma^2)}$;
                    \item objective: maximize NLL about $\bs{w}$ and $\sigma$;
                    \item solution: $\hat{\bs{w}}=\textbf{X}^{\dagger}\textbf{y}$ 
                    and $\hat{\sigma}^2=\dfrac{1}{N}\sum_{k=1}^{N}(y_k-\hat{\bs{w}}^{\top}\bs{x}_k)^2$
                \end{itemize}
                The quantity $\textbf{X}^{\dagger}$
                is pseudo inverse of matrix $\textbf{X}$, and if $\textbf{X}$ has full rank, it's value equals
                $(\textbf{X}^{\top}\textbf{X})^{-1}\textbf{X}^{\top}$.
                The matrix $\textbf{X}$ is called \textbf{design matrix}.
            \end{thm}
            Here, we can notice that the NLL is equal (up to unrelevant constants) to the RSS (residual sum of squares), given by
            \[
            \text{RSS}(\bs{w}) = \frac{1}{2} \sum_{k=1}^{N}(y_n - \bs{w}^{\top}\bs{x}_k)^2.
            \]

            If the output value is a vector, parameter estimation is done in a similiar way.

            \begin{thm}{MLE of Least Squares Linear Regression; vector output}{MLEofLSLR2}
                Let $\data=\{(\bs{x}_k, \bs{y}_k)|\ k = 1, 2, \ldots, N\}$, and let $n$ and $m$ be the dimension of $\bs{x}_k$ and $\bs{y}_k$.
                Let $\textbf{W} = (\bs{w}_1, \bs{w}_2, \ldots, \bs{w}_n)$ be \textbf{weight matrix}.
                \begin{itemize}
                    \item objective function:
                    \item objective:
                    \item solution:
                \end{itemize}
            \end{thm}
        
        Calculating the pseudo-inverse of $\textbf{X}$ could be established by directly calculating the inverse of $\textbf{X}^{\top}\textbf{X}$,
        but the given matrix may not have full rank.
        Instead, we have other methods that works well even when the $\textbf{X}^{\top}\textbf{X}$ is ill-formed.
        
        (W.I.P)
        
        \subsection{Further Informations}
            
            Least squares linear regression model has many variations, such as ridge regression, lasso regression, and robust regression, etc.
            It can represent and classify nonlinear data by using basis functions.
            These will be discussed soon.

    \section{Ridge Regression}

    \section{Lasso Regression}

    \section{Robust Linear Regression}

    \section{Bayesian Linear Regression}

    (W.I.P)

\chapter{Logistic Regression Models}

    \section{Binary Logistic Regression}
    \section{Multinomial Logistic Regression}
    \section{Robust Logistic Regression}
    \section{Bayesian Logistic Regression}

    (W.I.P)

\chapter{General Linear Models}

\chapter{Deep Neural Network}
    In this chapter, we will discuss about neural networks.
    Neural network is tool rather than the model itself.
    It can be applied at various problems such as regression, classification, and even generation.
    There are two main ways to combine DNNs with probabilistic models.
    \begin{itemize}
        \item define nonlinear functions inside conditional distributions;
        \item approximate the posterior distribution;
    \end{itemize}
    This is possible due to the \textbf{universial approximation theorem}, which says that
    neural network can approximate any given function.
    

    \section{Multilayer Neural Network}
    
        \subsection{Introduction}
        Here, the simplest form of neural network is given.

        \begin{defn}{Multilayer Neural Network}{NN}
            Let $\bs{x}$ be an $n$-dimensional vector input, and let $\bs{y}$ be an $m$-dimensional output vector.
            Let \{$\textbf{W}^{(k)}|\ k=1,2,\ldots,l$\} be weight matrices, let $h$ be an activation function,
            and let $f$ be an output activation function, whoose output is an $m$-dimensional vector.
            The following is the form of multilayer neural network model;

            \[
            \bs{y}(\bs{x},\textbf{W})=f(\textbf W^{(l)}(h(\textbf W^{(l-1)}\ldots h(\textbf W^{(1)}\bs{x}))))
            \]

            where $f(\cdot)$ and $h(\cdot)$ are evaluated on each vector element separately.
        
        \end{defn}
        Using this neural network, for example, in a regression problem, assuming that output value follows normal distribution,
        the model will look like this:
        \[
        p(t|\bs{x}, \textbf{W}) = \mathcal N (t|\bs{y}, \sigma^2)
        \]

        The neural network  can be easily recognized with following illustration: (PUT PICTURE HERE)

        \subsection{Backpropagation}

        \subsection{Further Informations}
        \subsubsection{List of Activation Functions}
            \begin{center}
                \begin{tabularx}{\textwidth}{XX}
                \toprule
                \textbf{Name} & \textbf{Definition}
                \\
                \midrule
                logistic sigmoid & $\sigma(a)=\dfrac{1}{1+\exp{(-a)}}$
                \\
                \midrule
                tanh & $\text{tanh}(a)=\dfrac{e^a-e^{-a}}{e^a+e^{-a}}$
                \\
                \midrule
                hard tanh & $h(a)=\max{(-1, \min{(1,a)})}$
                \\
                \midrule
                softplus (soft ReLU) & $h(a)=\ln{(1+\exp{(a)})}$
                \\
                \midrule
                ReLU & $h(a)=\max{(0,a)}$
                \\
                \midrule
                leaky ReLU & $h(a)=\max{(0,a)} + \alpha \min{(0,a)}$, where $0<\alpha<1$
                \\
                \midrule
                absolute & $h(a)=\vert a \vert$
                \\
                \bottomrule
                \end{tabularx}
            \end{center}

            ReLU (rectified linear unit) is the most widely used activation function, since it avoids vanishing gradient problem.
% move this subsubsection to next section.
        \subsubsection{Learning Methods} 
            \textbf{Representation Learning}

            \textbf{Transfer Learning}

            \textbf{Contrastive Learning}
    
        \subsubsection{Mixture Density Networks}
        We have been assuming that the output value follows distributions is Gaussian, which is unimodal.
        But, in real life, this may not match well.
        We can use \textbf{mixture models} to solve this problem. It is not a hard concept.
        What we have to do is just switch the output distribution from Gaussian to more complex, flexible distribution,
        such as \textbf{mixture Gaussian model}, which is multimodal.

    \section{Building Blocks of Neural Networks}


\end{document}